\chapter[Tang--Song Regional Networks]{Why Tang--Song China Became a Centre of Interregional Networks}
\section{A Porcelain Bowl from the Seabed}
First pick up an object: a porcelain bowl.

Small and open-mouthed, it carries brown and green patches and lively, loosely painted designs on blue-grey glaze. It is not precious; its surface is even somewhat rough. Not imperial tribute but a mass-produced commodity, it was fired near Changsha in Hunan over a millennium ago, thousands or tens of thousands emerging in each batch.

Its findspot records its journey. In 1998, local divers discovered a wreck near Indonesia's Belitung Island. Recovery revealed an Arab-style sailing vessel whose planks were sewn and bound with coconut-fibre rope instead of iron nails, a trading ship of the Persian Gulf and Indian Ocean. Its hold contained over sixty thousand neatly stacked ceramics, mostly Changsha ware. One bowl bore an ink inscription: ``Sixteenth day of the seventh month, second year of Baoli.'' Baoli was Tang Emperor Jingzong's reign title: the year was 826.

Trace the journey mentally. Inland Changsha lies beyond a thousand kilometres of mountains and rivers from the sea. A Hunan bowl first travelled south along the Xiang River system, eventually entered the Yangtze, and reached a port such as Yangzhou or Guangzhou. Loaded aboard a foreign merchantman, it sailed into the South China Sea and through waters near present-day Singapore. Thousands of kilometres from home, the ship sank, leaving it underwater for over eleven hundred years.

Was a cheap everyday bowl worth the distance and danger for a foreign trading fleet? Yes, and more than one ship made the journey. The wreck shows that ninth-century Chinese ceramics were already bulk commodities crossing half the globe, supported by invisible kilns, rivers, ports, taxes, markets, navigation, and thousands of livelihoods.

Our question lies within the bowl: how did Tang and Song China become the centre of the world's largest interregional network? Was that position natural or constructed? If constructed, which parts built it?

\section{High Tide at Quanzhou}
Now enter a scene three centuries later, moving from Changsha to the southeastern coast.

Time: a year in the twelfth-century Southern Song, the fifth lunar month, when summer monsoons begin. Place: Quanzhou, then called Zayton, or Citong, after the thorny, red-flowering coral trees planted throughout the city.

Before dawn, the tide surges into the Jin River estuary, lifting moored vessels almost a storey. The docks have been awake since midnight. Barefoot porters cross gangplanks with sacks of pepper, its sharp smell making newcomers sneeze. Nearby, wooden boxes hold frankincense and myrrh; sweet medicinal resin mixes with seawater and tung oil from the planks in the humid heat.

A Maritime Trade Office clerk stands in the yard with a ledger, followed by assistants carrying abacuses. His office resembles customs and a trade bureau combined: registering vessels, inspecting cargo, and extracting a proportion as tax in kind. Spice levies can bring considerable imperial revenue. Too many ships and calculations irritate him, but he knows half his salary comes from this pepper and frankincense.

The shipowner comes from Dashi, the Chinese name for the Arab regions. His halting Chinese does not hinder business. His surname is Pu, common among settled Arab and Persian descendants in Quanzhou, supposedly derived from Abu. This voyage brought him north from Champa, today's central and southern Vietnam. He sells spices inland and loads Jingdezhen qingbai porcelain, Zhejiang silk, and Fujian ironware, waiting for winter's northeast monsoon to return. A successful round trip takes two or three years; an unsuccessful one loses the ship to reefs, typhoons, or pirates. His profit is payment for risking life, not merely a price difference.

Towards town, the scene grows more mixed. A mosque's minaret rises among houses. Quanzhou still retains the Northern Song Qingjing Mosque's remains, among China's oldest surviving Islamic sites. Hindu temple carvings, Buddhist pagodas, and Daoist temples crowd one street. Outside the city, Mount Jiuri preserves Song inscriptions recording prayers for winds. Each summer and winter, local officials led shipowners uphill to worship: summer winds brought foreign vessels home, winter winds carried loaded ships out. Even government acknowledged the monsoon's hold on half the city's fortunes.

Remember this port. How did a city where officials, merchants, sailors, porters, and monks depended on one another grow? Why here in China's southeastern corner rather than elsewhere? The rest of our chapter answers its questions.

\section{What Makes a Centre?}
State the conflict before rushing to answers.

A convenient answer says China had always been powerful, Tang and Song marked an ancient peak, and centrality followed naturally. Inspiring, perhaps, but vulnerable to one question: during Quanzhou's greatest prosperity, the Song was notably weak militarily.

Consider established facts. The Northern Song never recovered the Sixteen Prefectures of Yan and Yun, facing Liao to the north and Western Xia to the northwest. The 1005 Chanyuan agreement exchanged annual Song silver and silk payments to Liao for frontier peace. In 1127, Jin troops captured Kaifeng and two emperors, ending the Northern Song. The surviving court fled south and settled at Lin'an, today's Hangzhou, ruling the south for over a century. This is the regime behind our admired Quanzhou, barely able to defend half the country.

The Tang presents another tension. Chang'an's splendour was real, but the later dynasty was troubled. The An Lushan rebellion of 755--763 broke the empire across its middle. For over a century afterwards, central rulers struggled with autonomous regional military governors, their revenue at times hanging on wealthy southeastern regions and Grand Canal grain transport.

Thus militarily battered Song China nevertheless reached the imperial era's greatest heights in economic scale, cities, and overseas trade. Chang'an, Kaifeng, Hangzhou, Quanzhou, and Guangzhou held hundreds of thousands, sometimes over a million, while contemporary Europe's largest cities often had only fractions of those populations.

At least one conclusion follows: power bundles several different things. Military strength, economic size, institutional capacity, and attraction to neighbours---drawing people to trade, learn your writing, or worship at your temples---are distinct. They do not always rise or fall together.

Our real question is therefore: \textbf{what components establish a network centre?} Geography, institutions, market demand, or their combination? What weight does each carry? This also concerns today's financial centres, container ports, and resource-rich countries remaining on network margins.

Before proceeding, place a bet: if geography, institutions, or luck must be the most important component, which do you choose? Remember it and check again at the end.

\section{Court and Merchants: Who Supports Whom?}
To understand centres, listen first to contemporary arguments. At least three different voices within Tang--Song China left traces concerning commerce and foreign trade.

\textbf{Voice one: scholar-officials' suspicion.}

For a long time, central dynasties ranked scholars, farmers, artisans, and merchants as four social orders, merchants last. Official policy promoted agriculture and restrained commerce: merchants could not wear fine clothes, use fine carriages, or easily see their descendants enter office.

Do not immediately mock this. \textbf{State their case in full}, since equating restraints on commerce with ancient stupidity is unfair.

Agrarian advocates had at least three reasons. First, food sustained life. With low yields and human or animal transport, a poor harvest could become famine, displacement, and rebellion. Tying as many people as possible to grain production was crude but reliable insurance. Second, taxation was practical: households and adult males fixed on land were easy to register and assess; merchants moving between Guangzhou and Yangzhou with money in chests were not. Third, fairness mattered. A farmer might save only a few measures of rice after a year's labour, while one spice deal earned a merchant ten years of that income. Without restraint, why would farmers stay? If everyone traded, who would grow food?

This was rational risk management under agrarian constraints, not foolishness. Its fatal blind spot was treating farming and commerce as rivals rather than mutually sustaining. Merchants brought grain to disaster areas; southern weaving and transport clothed northern troops. The court increasingly found commercial revenue, especially Maritime Trade Office levies on spices and ceramics, indispensable. Song policy's relatively loose restraints and diligent taxation followed an accounting realisation, not sudden enlightenment.

\textbf{Voice two: merchants and sailors.}

They left objects rather than official memorials: wind-prayer inscriptions, pepper heaps, Pu-family houses, and Arab travellers' reports.

In the mid-ninth century, an Arab named Sulayman recorded widely circulated eastern observations, later commonly called Accounts of China and India. He described Guangzhou's foreign merchant population and official appointment of a Muslim community leader to administer internal affairs by its own laws and customs. He marvelled at porcelain so thin that water's colour showed through its walls. Here Chinese government appears shrewd rather than exclusionary: welcome newcomers, allow autonomy, collect taxes. It is the same face as Quanzhou's Maritime Trade Office.

\textbf{Voice three: unnamed kiln and ship workers.}

The network rested on countless ordinary people absent from histories. Their voices also survive, written on ceramics. Changsha wares bear popular verses painted in brown, including this poem:

\begin{quote}
When you were born, I was not; when I was born, you were old. You regret my late birth; I regret your early birth.
\end{quote}
Inscribed on a Tang ceramic ewer, it has no known author, perhaps a kiln worker or a customer's casually requested love verse. Consider its significance: a cheap commodity destined for half a world away carries ordinary lovers' regret. Maker, writer, buyer, and the sailor who sank off Belitung all belonged to the same network. They may not have known grand terms such as empire and network, but every thread passed through their hands.

Together the voices complicate the image of a supreme celestial empire. Government distrusted merchants yet depended on their taxes; merchants occupied low social rank but held real wealth; ordinary people bore risks and shared prosperity. Centrality arose not by command but through these forces pulling against and using one another until a balance emerged.

\section[Thinking Bridge: Invention and Diffusion]{Thinking Bridge: Separate Invention from Diffusion}
Discussions of Tang--Song advancement love invention lists: paper, printing, gunpowder, compasses, paper money. As the list lengthens, explanation grows muddier. Here is a portable cure for list-making disease.

For any technology, separate one question into three:

\begin{itemize}
\item \textbf{Level one: who first invented it?}
\item \textbf{Level two: where and why was it widely adopted?}
\item \textbf{Level three: what changed after adoption?}
\end{itemize}
Answers often point to different places, with levels two and three more important than one. Invention is a point, diffusion a line, widespread application the surface that changes the world. Judge technological capacity by diffusion speed and application depth, not counts of firsts.

Apply this to \textbf{an easily misjudged counterexample: movable type}.

The invention list puts Chinese movable type in the eleventh century. Shen Kuo's Dream Pool Essays records: ``During Qingli, a commoner named Bi Sheng also made movable plates.'' Qingli fell in the 1040s, roughly four centuries before Gutenberg. If earlier invention meant leadership, China should have led the printing revolution throughout.

Yet until the nineteenth century, \textbf{woodblock printing} dominated China, while Bi Sheng's clay type and later wooden type remained secondary. The explanation is not stupidity but level two. Tens of thousands of Chinese characters required vast type stocks and literate compositors, creating formidable costs. Woodblocks suited long-selling classics, calendars, and medical books: carve once, store, print as needed for decades. To publishers, blocks were dependable fixed assets and movable type a risky gamble. Markets calculated costs, not the prestige of advancement.

Gutenberg encountered different conditions in fifteenth-century Europe: only dozens of letters, little type, lower composition barriers, and suddenly favourable economics. Printing spread like wildfire, producing more books in fifty years than the previous millennium's scribes. Korea took another route: thirteenth-century Goryeo adopted metal type because government financed Buddhist and classical publication, without needing market approval.

Three inventions, three fortunes, determined not by dates but by language, market structure, and institutions: \textbf{levels two and three}.

Apply the tool again to the compass. Near the end of the Northern Song, Zhu Yu's Pingzhou Table Talks recorded Guangzhou navigators:

\begin{quote}
Shipmasters know geography: by night they observe stars, by day the sun, and in darkness or cloud they consult the south-pointing needle.
\end{quote}
This is among the earliest clear accounts of navigational compass use. At level one, China long knew magnetism's directional property. Levels two and three changed the world: Song ocean-going traders adopted all-weather navigation, then Arab trading routes carried it to the Mediterranean, making it crucial to later European ocean voyages. Seventeenth-century English philosopher Francis Bacon remarked, broadly, that printing, gunpowder, and the magnetic needle transformed the world though their origins were obscure by his time. He was right: all three originated within this eastern network, with China a major node.

Keep the tool ready. Evidence and explanations ahead will repeatedly require it.

\section{The Grand Canal in a Poem}
Now read a primary source. Late-Tang poet Pi Rixiu's ``Reflecting on Antiquity beside the Bian River'' concerns the Grand Canal:

\begin{quote}
All say this river brought the Sui down; its flowing waters still sustain a thousand miles. Without the water palaces and dragon boats, comparing the achievement with Yu's would not be excessive.
\end{quote}
Take it line by line. Everyone says the canal destroyed the Sui: a prevailing contemporary interpretation supported by evidence. Emperor Yang conscripted millions to dig it, ruining countless families. Labour demands and wars against Goguryeo together crushed the dynasty little more than a decade after construction began.

Yet its waters still sustain a thousand miles. Writing two centuries after the Sui fell, Pi knew Tang finances depended on it: southern grain, silk, and taxes travelled north by ship, feeding Chang'an's court and frontier armies.

The final lines give Pi's own judgement: without pleasure voyages to Yangzhou in dragon boats, Emperor Yang's canal-building achievement could fairly be compared to Great Yu's flood control.

Notice the unusual distinction. What happened: conscription, rapid dynastic collapse, continued canal flow. Contemporary understanding: evidence of tyranny. Evaluation: Pi calls the engineering meritorious and pleasure-seeking culpable. One work can be both indictment and artery, depending on its aspect and how long after completion you live.

This hardware supplies our first puzzle piece. The canal joined the Yellow and Yangtze rivers and southern waterways into an inland transport network carrying grain, ceramics, salt, iron, people, and documents. It brought Changsha bowls to ports and Quanzhou spices to Kaifeng markets. Northern Song Kaifeng placed its capital at the canal hub, a million-person metropolis fed by logistics, unique in the contemporary world.

Add an outsider's testimony. Late-thirteenth-century Venetian merchant Marco Polo visited Quanzhou and praised Zayton's port, roughly saying that for every shipload of pepper reaching Alexandria, Christendom's greatest port, a hundred or more reached Zayton. Discount travellers' numbers: the hundredfold comparison is not literal. Yet the impression of scale is credible: Quanzhou's trade dwarfed the Mediterranean world in his eyes.

\section{Three Explanations, Each with Gaps}
With evidence assembled, ask again why Tang--Song China became an interregional network centre. These three explanations differ genuinely, each holding part of the truth and falling short elsewhere.

\textbf{Explanation one: geography and waterways.}

Two great east--west rivers with plains between became highways once joined by canals. An indented southeastern coast offered good harbours facing the busy Indian Ocean--South China Sea route. Reliable seasonal monsoons supplied free engines. Geography explains why this network grew here rather than in deserts or highlands.

Its limit: geography was already there. Qin and Han China had both rivers too. Why the Tang--Song surge? Ninth- and thirteenth-century monsoons were much the same. Geography sets the stage, not the script: it explains possibility, not timing.

\textbf{Explanation two: institutional work.}

Tang--Song institutions offered substantial components. The canal was national infrastructure. Examinations, established under Sui and Tang and matured under Song, drew regional elites into one system, with anonymisation and recopying to reduce cheating: names sealed, scripts recopied before marking, opening routes for poorer candidates. This bound the empire and generated literacy. Maritime Trade Offices regulated and fostered commerce. Early Northern Song Sichuan merchants, burdened by iron cash, jointly issued paper deposit certificates, jiaozi. In 1023 government took over and established the Yizhou Jiaozi Office, issuing what is widely recognised as the earliest official paper currency. Woodblock printing sharply cut book prices, making knowledge affordable to ordinary people for the first time.

Its limit: institutions are not automatic ticket machines. Excess paper issues under the later Southern Song and Yuan extracted wealth and destroyed credit. Maritime Trade Offices could close when courts turned restrictive. There is also circularity: good institutions explain prosperity, which proves the institutions good. Where did institutions come from? Why did Song rulers loosen commercial restraints? This explanation cannot answer.

\textbf{Explanation three: demand and networks.}

Zoom out: prosperity was not a solo performance. From the eighth through thirteenth centuries, both ends of the Indian Ocean--South China Sea routes flourished. Abbasid Baghdad formed another pole; Persian Gulf, western Indian, and Southeast Asian ports prospered. Eurasia needed Chinese ceramics, silk, and ironware, while China needed spices, medicines, pearls, and silver. China did not simply radiate influence outward; an interregional network matched mutual demands, with China its largest supply node.

Its limit: demand existed everywhere, so why could China meet it? Industrial, standardised production of sixty thousand bowls for one ship required organised kilns, rivers, and ports. This returns to institutions. Worldwide demand alone does not explain Chinese supply.

These explanations are three storeys. Geography determines where centrality is \textbf{possible}; institutions determine whether possibility is realised; network demand determines its scale. Remove one and the building fails. This differs from the empty phrase ``many factors combined'': separating levels allows each to be tested independently.

\section[World-System Centres]{Further Challenge: What Is a World-System Centre?}
Now introduce the chapter's theory, from twentieth-century history and sociology: world-systems analysis.

American scholar Immanuel Wallerstein argued in the 1970s that modern history cannot be divided country by country, because economies form an interlocking system. It contains a \textbf{core}, performing profitable, complex work with strong bargaining power; a \textbf{periphery}, supplying raw materials and cheap labour with weak bargaining power; and a \textbf{semi-periphery} between. What matters is not national industrious character but systemic position, determining a country's share of the network's gains.

Was such an interregional system invented by modern Europe? No. Janet Abu-Lughod's Before European Hegemony looks to the thirteenth century: before Europe's rise, interlocking Eurasian circuits already connected northwestern Europe, the Mediterranean, the Indian Ocean, and China, its weightiest node. Her strongest point is that \textbf{this system had no hegemon}. Later European dominance grew after the old network loosened. Centrality is neither natural nor destined, only a temporary position.

Applied to Tang--Song China, this immediately does three useful things.

First, it corrects our wording. More accurately than ``China was the centre,'' China was one of the largest, most productive nodes in a multicentred Eurasian network, alongside Baghdad, Cairo, and western Indian ports. This dismantles the one-way celestial-empire picture. Quanzhou's Pu merchants were trading partners, not tributary servants; networks ran both ways.

Second, it explains simultaneous military weakness and prosperity. Interstate battle results and network positions based on supply, market size, and institutional costs belonged to different contests. Battlefield losses did not directly subtract network points.

Third, it demystifies centrality. A centre is a position, not a trophy. Routes shift, credit and monetary rules change, other nodes grow, and centres move. Tang--Song centrality was one stop in a long movement, not its endpoint.

Weigh the theory's limits too. It draws great arteries well but can reduce people inside them to factor flows: bowls become commodities, kiln workers labour, wrecks trade losses. The love-poem ewer reminds us that every link contains named people who fall in love and drown. Theory is a map, not the territory: it explains centre formation, not whether prosperity treats people fairly. That judgement is yours.

\section{Scriptures and Images Travelled Too}
The network carried more than goods. Ideas followed the same sea routes.

Tang thought shared a stage among three traditions. Confucianism supplied official learning and exams. Buddhism, arriving from India under the Han, developed deep Chinese roots and forms such as Chan; Xuanzang's Indian journey and major translation work belonged to the seventh century. The Tang's Li rulers, claiming descent from Laozi, strongly promoted native Daoism. The traditions argued and borrowed vocabulary without eliminating one another: normal internal diversity, not a monolith.

Ideas travelled with people. Japanese missions repeatedly sent students and monks who spent a decade or more in Chang'an and returned with scriptures, Chinese books, calendars, and capital plans. Heijō-kyō and Heian-kyō's grid streets echo Chang'an. Eighth-century monk Jianzhen attempted six crossings from Yangzhou, persisting despite blindness before bringing Buddhist discipline and Tang culture to Japan. Korean scholars came to study and sit examinations; some passed, held Tang office, and returned with their learning.

Under Song, sustained Buddhist and Daoist challenges prompted Confucian reconstruction. Northern Song brothers Cheng Hao and Cheng Yi reinterpreted classics; Southern Song Zhu Xi synthesised what later became Neo-Confucianism, or the Learning of Principle. Rather than strange supernatural matters, it discussed li and qi, principle and material force, mind and cultivation, making a system capable of confronting Buddhist and Daoist thought. This was competitive renewal, not mere revival: a tradition's strongest voice often follows its greatest challenge.

The learning soon travelled old routes overseas. Zhu Xi's teachings became official Joseon Korean learning, shaped Vietnamese education and exams, and gained prominence in shogunal Japan. Examination classics and scholarly political vocabulary became transnational academic currency, circulating like ceramics and coins. Here is the other side of the outline's claim that innovation depends on knowledge transmission: \textbf{institutions, markets, and knowledge are three cargoes on one network}. Exams generated readers, printing lowered prices, cheaper books enlarged readership, and readers supported exams and printing. The links were inseparable.

Add relations between scholar politics and ordinary people. Examination-trained Song officials had enough political influence to openly oppose imperial reforms. Leaders on both sides of the Northern Song's great reform debate came through regular examinations, contesting our central question: should government intervene more deeply in markets? One camp supported state finance, lending, and regulation; the other feared competition with the people and harm to society. No true victor emerged, but Song thinkers were already seriously debating strikingly modern network-governance problems: openness or control, markets or planning.

Ordinary people appear elsewhere than memorials: boatmen's chants on the Bian, tea and drink stalls in Kaifeng's night markets, names beneath Quanzhou's wind inscriptions, and poems on export ewers. Prosperity's weight was unequal. Customs revenue, kiln orders, and dock jobs were real; so were conscription, shipwrecks, and inflation. Remember that inequality for our final question.

\section{Today: Ports, Exams, and QR Codes}
This thousand-year-old question lies beneath your feet.

Consider ports. In 2021, ``Quanzhou: Emporium of the World in Song--Yuan China'' entered the World Heritage List, including the Maritime Trade Office, Qingjing Mosque, and Mount Jiuri wind-prayer sites. Hundreds of kilometres away, Shanghai has led world container throughput for years, followed by Ningbo-Zhoushan and Shenzhen. A network once fronted by Quanzhou and Guangzhou now covers the southeastern coast at hundreds or thousands of times the scale. Night-time cranes and streams of trucks are twenty-first-century wind-prayer inscriptions, appealing now to supply-chain punctuality rather than gods.

Consider money. Northern Song jiaozi began the earliest paper-currency experiment, later undermined by Yuan and Ming overissue and credit collapse, leading to paper's temporary retreat. A millennium later, paper money's birthplace has become a leader in mobile-payment penetration, with QR codes at street stalls and difficulty finding cash change. Our tool reveals a reversal: China invented neither the internet nor smartphones, yet led in mobile payments' \textbf{diffusion speed and application depth}. Level one's invention answer is America; levels two and three, adoption and consequences, point here. The old pattern persists: application networks, not invention lists, transform the world.

Consider exams. Imperial examinations are gone, but institutional genes remain: mass selection through standardised testing and faith in upward routes for poor children. Echoes are not identity: classical essays for official appointment differ enormously from today's subjects and destinations. Yet debates over fairness and a single decisive test revisit Song questions. Standardisation offers fairness and rigidity; upward routes must exist, but who shapes them?

Finally, consider knowledge. Woodblocks cut prices and expanded readership; the internet cuts transmission costs further. Once classics were unaffordable; now information exceeds reading capacity and truth is hard to distinguish. The old network question changes clothes: with universal connection, scarcity shifts from access to discernment. This is precisely the capacity the whole book seeks to develop.

\section[Your Question: Pricing Centrality]{Your Question: Who Prices Admission to the Centre?}
Return to the porcelain bowl.

An ordinary Hunan bowl, underwater in Indonesia for eleven hundred years, carries a whole network: canals, ports, taxes, monsoons, an Arab captain, Changsha workers, an unnamed poet. We have examined its construction: geography below, institutions as beams, demand above, ideas travelling aboard ships. A centre is no civilisational personality or eternal throne, but a temporary position maintained by countless components.

Now answer a question with no standard response:

\textbf{If becoming a network centre brings greater openness and opportunity but also dependence, risk, and competition, should a society pursue that position?}

Before deciding, weigh every side again:

You hold reasons in favour: prosperity, jobs, influence, and choice. Quanzhou porters and Changsha workers earned their food through the network. Song commercial loosening shows openness's potentially enormous returns. World-systems theory also reminds you that positions are scarce: others occupy those you leave vacant.

You also hold cautions. Centrality is never free. Militarily weak Song paid annual transfers and war bills. Currency collapse cost ordinary savers. Denser networks transmit failures faster: one wreck takes cargo owners, workers, and lenders down together. The poem-bearing ewer reminds you that some prosper greatly while others receive only the lament about lovers born too far apart.

A third option is easily overlooked: perhaps the question is mistaken. If a centre is a moving position rather than a trophy, might pursuing permanent centrality itself be an illusion, like rowing endlessly upstream?

What is your answer, and why? Identify your reasoning's level: geography, institutions, or network demand. Making clear where you stand is itself this chapter's central lesson.
